% =====================================================================
%  第 5 章  评价方案
% =====================================================================

\section{评价方案}
\label{sec:evaluation}

第\ref{sec:results}节报告了确定性骨架的结果与一次完整的智能体运行。本章说明如何判断完整系统
（含大语言模型）运行结果的好坏。消融实验与筛选环节的安慰剂对照均已完成，
人为注入假命题的检验留待后续。

\subsection{评价指标的选取}

最直观的指标是系统提出了多少条后来被证明的猜想。该指标在课程周期内无法评估，因为一条猜想从提出到获证通常需要数年。

退一步，可以统计系统提出了多少条新猜想。该指标同样不可靠：新颖性需要查证文献，而“未见记载”往往只意味着检索不充分。更严重的是，该指标会诱导系统追求冷僻而难以验证的形式，因为形式越冷僻越显得新颖。

因此本课题采用一个较为保守但可操作的标准：\textbf{系统提出的猜想，能否使一位具备图论基础的研究者在不到一小时内判断其真伪}。该标准同时排除两类失败：命题过弱时，实为已知结论的改写，不构成发现；命题过怪时，涉及无人研究过的量，无从下手，亦不构成有效产出。据此，一条合格的猜想应满足：形式简洁，所涉不变量均为常见量，在现有数据上高度可信，且已有反例搜索未能将其推翻。

此外，系统须通过一项正确性检验：在正式运行之前，先令其“发现”已知成立的结论，
例如命题 \ref{prop:chi-alpha}。该项仅用于确认流程无误，不作为成果。
第\ref{sec:results}节已在不调用模型的基线上完成此项检验，
五条已知结论均被复现。
若系统无法找到这些不等式，或找到后又被自身推翻，则说明实现存在缺陷，
后续结果均不可信；反之，系统能重新发现这些结论，
也仅说明其工作正常，不足以说明其有效。

\subsection{对照实验}

判断一条候选是否“无效”需要标准答案，而真实数据上没有标准答案。
本课题换一个角度：构造一份安慰剂数据。
把目标列（色数）在整份数据上随机置换之后，色数与其余八个不变量之间
不再有任何关联，任何在这份数据上“成立”的候选都必然是假发现。
于是安慰剂数据上保留的条数，直接就是假发现的条数。

在同一流程上分别统计三种判定方式保留的条数：
不作校正（$p \leq 0.05$ 即保留）、BH 校正取 $q = 0.1$、BH 校正取 $q = 0.2$。
三组种子下的结果如表\ref{tab:placebo} 所示。

\begin{table}[H]
\centering
\small
\caption{安慰剂数据与真实数据上的保留条数}
\label{tab:placebo}
\begin{tabular}{@{}lcccccc@{}}
\toprule
判定方式 & \multicolumn{4}{c}{安慰剂（目标列随机置换）} & 真实数据 \\
\cmidrule(lr){2-5}
 & 种子 0 & 种子 1 & 种子 2 & 均值 & \\
\midrule
进入筛选的候选 & 27 & 0 & 22 & 16.33 & 17 \\
不作校正（$p \leq 0.05$） & 0 & 0 & 0 & 0.00 & 7 \\
BH 校正（$q = 0.1$） & 0 & 0 & 0 & 0.00 & 7 \\
BH 校正（$q = 0.2$） & 0 & 0 & 0 & 0.00 & 7 \\
\bottomrule
\end{tabular}
\end{table}

两点结论。其一，置换检验确实在起作用：
安慰剂数据上一条也没有保留，真实数据上保留了 7 条。
这说明该环节不是形式上的装饰。

其二，本次实验未能证明多重检验校正本身的必要性。
在真实数据上，不作校正与 BH 校正的结果完全相同（都是 7 条）。
原因在于置换次数 $B$ 决定 $p$ 值能取到的最小值 $1/(B+1)$：
$B = 200$ 时下限为 0.005，而 BH 要求存在 $i$ 使 $p_{(i)} \leq i q / m$，
故只有当候选数 $m \leq (B+1)q$ 时才可能保留任何候选。
$q = 0.1$ 时这一上限是 20，本次 $m = 17$ 恰好落在边界内，
且 7 条候选的 $p$ 值都取到了下限，校正因而没有产生任何差别。

这一发现是实验的产物，也构成对方法的一项限制：
若候选数超过约 20 条，或置换次数不足以让 $p$ 值下探到 $i q / m$ 以下，
BH 会变得极其保守，甚至一条也不保留。
换言之，该环节的有效性依赖于置换次数与候选规模的匹配，
而这在事先并不显然。后续工作将把 $B$ 提高到 2000 以上，
并在一批人为注入的假命题上检验校正的筛除能力。

\subsection{拟报告的负面结果}

本课题拟明确报告以下几类负面结果，而不加掩饰。

\begin{itemize}
  \item 大语言模型的贡献有限。消融实验（表\ref{tab:ablation}）显示，
  剔除恒真式之后模型提议相对随机抽样未见可测量的优势，
  这一结论已在第\ref{sec:results}节如实报告。
  \item 统计筛选的代价过高。安慰剂实验（表\ref{tab:placebo}）显示，
  置换次数与候选规模若不匹配，BH 会变得极其保守，甚至一条也不保留。
  本次实验亦未能证明校正相对不校正的必要性，这一局限已如实记录。
  \item 系统仅在小规模图上有效。若全部结论均局限于七、八阶图，其数学意义有限，这一点须予说明。
\end{itemize}

上述结果并不构成本课题的失败。一份能够明确说明方法在何种范围内有效、在何种范围内失效的报告，较之仅展示成功案例的报告更具参考价值。
